\section{Marco teórico}

El análisis financiero es una lectura técnica de los estados financieros que busca explicar la situación de una empresa y apoyar decisiones económicas. El Marco Conceptual NIIF sostiene que la información financiera es útil cuando sirve a los usuarios para tomar decisiones.

La NIC 1 destaca la importancia de la presentación y comparabilidad de los estados financieros. Esta comparabilidad permite analizar cambios entre periodos y estudiar la estructura interna de los estados.

En el contexto peruano, el Plan Contable General Empresarial organiza cuentas contables que sirven como base para preparar estados financieros. Por ello, existe una relación entre registro contable, clasificación de cuentas, estados financieros y análisis.

Fuentes como Amat, Edwards, Hermanson y Maher, y la bibliografía contable registrada en el proyecto resaltan que el balance y el estado de resultados son instrumentos centrales para comprender la posición financiera, el desempeño y la utilidad empresarial.
